%%%%%%
%
% $Autor: Wings $
% $Datum: today $
% $Directory: KeywordSpottingData $
% $Version: 2 $
% $Autor:  $
% $Review date: $
%
%%%%%%

\chapter{Keyword Spotting Data}

The data used for keyword spotting in this project is primarily audio samples collected from various sources. The audio signals are pre-processed and transformed into a suitable format for machine learning models. These signals contain both command keywords and background noise to improve model robustness.

\section{Origin}

The dataset used in this project originates from publicly available keyword spotting databases such as Google Speech Commands Dataset. These datasets provide labeled audio samples for various spoken keywords, background noise, and silence.



\section{Structure}

The data structure consists of labeled audio files representing specific keywords (e.g., ``yes'', ``no'', ``stop'') and non-keyword background sounds. Each sample is represented as a spectrogram or \ac{MFCC}, depending on the preprocessing pipeline used. Labels are stored in a structured format such as CSV or JSON.

\medskip

Key features extracted from the audio samples include:

\begin{itemize}
	\item \textbf{\ac{MFCC}:} Used to capture spectral properties of the audio.
	\item \textbf{Spectrogram data:} Visual representation of the signal's frequency content over time.
	\item \textbf{Zero-crossing rate:} Measures signal oscillation rate.
	\item \textbf{Energy and entropy of energy:} Highlights amplitude variations.
\end{itemize}

\section{Size}

The dataset typically includes:

\begin{itemize}
	\item Audio Samples: 1-second clips recorded at a sampling rate of 16 kHz.
	\item Dataset Size: Approximately 50 MB to 100 MB for small-scale training and up to several GBs for more extensive datasets.
	\item Keywords: 10 to 20 classes, including a background noise class.
\end{itemize}


\section{Format}

The data is stored in the following formats:

\begin{itemize}
	\item Audio Files: \texttt{.wav}, \texttt{.mp3}
	\item Labels: \texttt{.csv}, \texttt{.json}
	\item Processed Features: \texttt{.npy} (NumPy arrays) or \texttt{.tfrecord} (TensorFlow format) for integration with machine learning pipelines.
\end{itemize}

\section{Data Types}

The dataset includes:
\begin{itemize}
	\item \textbf{Audio recordings:} 16-bit PCM WAV files sampled at 16 kHz.
	\item \textbf{Labels:} Textual labels indicating the keyword spoken.
	\item \textbf{Metadata:} Information such as recording environment and speaker demographics.
\end{itemize}


\subsubsection{Quality}
\begin{itemize}
	\item \textbf{High-quality recordings:} Majority of audio samples are clear with minimal noise.
	\item \textbf{Potential inconsistencies:} Some recordings may have background noise or low volume.
	\item \textbf{Annotation accuracy:} Labels have been verified, ensuring high confidence in correctness.
\end{itemize}

\subsubsection{Quantity}
The dataset consists of:
\begin{itemize}
	\item \textbf{Number of samples:} Over 65,000 audio recordings.
	\item \textbf{Class distribution:} Balanced distribution of keywords, silence, and background noise categories.
\end{itemize}

\subsubsection{Fairness and Bias}
\begin{itemize}
	\item \textbf{Speaker diversity:} Dataset includes speakers of different ages, genders, and accents.
	\item \textbf{Potential bias:} Overrepresentation of certain demographics could lead to model biases.
	\item \textbf{Mitigation:} Augmentation techniques are applied to simulate diverse environments.
\end{itemize}

\section{Outliers and Anomalies}


The dataset may contain:

\begin{itemize}
	\item Noisy Samples: Background interference from environmental sounds.
	\item Missing Labels: Some files may lack corresponding annotations.
	\item Silent Samples: Audio clips without any recognizable sound due to recording errors.
\end{itemize}

\begin{itemize}
	\item \textbf{Outliers:} Samples with extreme background noise or unusual audio patterns were flagged.
	\item \textbf{Anomalies:} Detection algorithms were used to identify and either correct or exclude problematic samples.
\end{itemize}

Preprocessing steps, such as noise filtering and normalization, are applied to mitigate these anomalies.

\section{Origin}

The data is sourced from open datasets such as the Google Speech Commands Dataset and supplemented with custom recordings for specific application needs. The open datasets provide a rich variety of keywords and background noise samples, ensuring diverse and realistic training conditions.


